\documentclass[12pt]{article}
\usepackage[utf8]{inputenc}
\usepackage[T1]{fontenc}
\usepackage{hyperref}
\title{Klingon Language Resources for K-12 Curriculum Development}
\author{Compiled by Claude Sonnet 5 (claude-sonnet-5) for Andr\'as Kornai}
\date{}
\begin{document}
\maketitle

\section{Introduction}
Klingon (tlhIngan Hol, ISO 639-3 tlh) is a constructed language. Marc Okrand devised it for Star Trek. It has no community of native speakers. Learners and enthusiasts study it. The Klingon Language Institute (KLI) is the main organization.

\section{Primary Educational Resources}
\begin{itemize}
\item A University of Arizona community classroom syllabus for a Klingon course, by Warner: \url{https://communityclassroom.arizona.edu/sites/default/files/2023-07/KlingonSyllabus_Warner.pdf}
\item A KLI mailing-list thread about a Klingon course at Signum University: \url{https://lists.kli.org/archives/list/tlhingan-hol@lists.kli.org/thread/OHW6636JLX4GPE5STNQVV6Q4U5KSMXKF/}. The search did not show whether the course is still offered.
\end{itemize}

\section{Linguistic and Reference Materials}
KLI list posts treat Marc Okrand's dictionary as the essential reference for the grammar. I did not find a direct link to the book. The KLI list archive is at \url{https://kli.org/tlhIngan-Hol/1998/December/msg00093.html}.

\section{Religious and Community Resources}
The KLI's quarterly journal HolQeD is described in an early list post: \url{https://kli.org/tlhIngan-Hol/1993/July/msg00031.html}. The same post mentions a free postal course of eleven lessons. I found no current details on it.

\section{Digital and Audio Resources}
A 2022 course syllabus recommends the free app ``tlhIngan Hol boQwI'\,'' as a dictionary. The search gave no direct store link.

\section{Curriculum Development Context}
The language has no native speakers, so ``native'' materials are those of the constructed-language community. Most of what was found is old mailing-list material and one university syllabus. Current KLI pages would need to be checked.
\end{document}